\subsection{How many arithmetic tests may fail in a complete window}
\label{scope:sparse-window}

Fix the exponent, the retained family and the full page, Behrend
and Fourier hypotheses of Proposition~\ref{scope:fixed-powers}.
Put \(\beta=1/2+97\varepsilon/16\) and fix an integer \(r\ge1\).
In a complete window write \(Z=Y^{3/2}\) and
\(\mathcal N_Y=\{m\in\mathbf Z_{>0}:Y\le m^r\le Z\}\).
At \(x=m^r\) let \(Q(x)=\bar D(x)/E_x\) on the input mean hosts.
On the realized overlap hosts one may equally take
\(Q(x)=D_{u(x)}(x)/E_x\) for any selection of translates, including
a minimizing or maximizing translate chosen at each \(x\).
All estimates below are uniform in that choice.

\begin{lemma}[Detecting a signed excursion in a fixed multiplicative interval]
\label{scope:fixed-ratio}
There is a finite \(C=C_{\psi,\beta}>1\) with the following property.
If \(\Theta>\beta\), then for every \(\beta<\theta<\Theta\) and
all sufficiently large \(X\) there is a \(t\) in \([X,CX]\) with
\(F(t)>t^\theta\). One may keep both ends of this interval a fixed
positive logarithmic distance from \(t\). If the real parts of all
real zeros are strictly below \(\Theta\), the same assertion holds
with \(F(t)<-t^\theta\), at a possibly different point.
\end{lemma}
\begin{proof}
We use the positive measure construction in the project's
\emph{Zeta positive-part sieve}, \S\S2--3, for the full retained
Dirichlet family, taking account of real poles. Fix
\(\beta\le\sigma<\Theta\), initially \(\sigma=\beta\).
Among the distinct poles \(\rho=b_\rho+i\gamma_\rho\) of
Proposition~\ref{an:poles}, take those with \(b_\rho>\sigma\) and
\(\gamma_\rho>0\). Then \(L=\sum_\rho\pi/\gamma_\rho<\infty\).
Indeed, for the fixed finite family Montgomery's density theorem
bounds their count up to height \(T\ge2\) by
\(O_{\psi,\sigma}(T^{3(1-\sigma)/(2-\sigma)}\log^{K}(2T))\)
for an absolute \(K\).
The exponent is less than one for \(\sigma>1/2\), so partial
summation proves convergence. See
\href{https://link.springer.com/book/10.1007/BFb0060851}{H.~L. Montgomery,
\emph{Topics in Multiplicative Number Theory}, Chapter 12}.

Write \(z_\rho=\rho-\beta\), \(l_\rho=\pi/\gamma_\rho\). Take the
sum of independent variables with the following two-point laws:
\[
   \frac{e^{(b_\rho-\beta)l_\rho}\delta_0+\delta_{l_\rho}}
        {1+e^{(b_\rho-\beta)l_\rho}} .
\]
The resulting probability measure \(\nu\) is supported on \([0,L]\).
Its entire transform \(M(z)=\int e^{zv}\,d\nu(v)\) vanishes at
every \(z_\rho,\overline{z_\rho}\) on this list by independence,
and is strictly positive on the real axis. Convolve \(\nu\) with
uniform probability on \([-D,D]\), where \(D=L+1\). The density
\(B\) so obtained is supported on \([-D,L+D]\) and equals
\(1/(2D)\) almost everywhere on \([0,L]\).

For \(g(u)=e^{-\beta u}F(e^u)\), the Hadamard expansion in
Proposition~\ref{scope:fixed-powers} has an absolutely summable
coefficient for each nontrivial zero, given by \eqref{eq:residue}.
The constant, zero at the origin and negative trivial zeros contribute
\(O_\psi((1+u)e^{-\beta u})\) for \(u\ge1\). Consequently
\(\int B(v)g(u+v)\,dv\) is the sum of the terms of the real poles
to the right of \(\sigma\), each with strictly positive coefficient
and exponent \(\rho-\beta\), plus \(O(e^{(\sigma-\beta)u})\).
If a real pole has real part \(\Theta\), this proves the required
positive excursion by taking \(\sigma=\beta\).

Otherwise choose a nonreal \(\rho_*\) with real part \(b_*>\theta\)
larger than \(\sigma\) and the real parts of all real poles. Such a
choice exists since the real poles in \([\beta,1]\) are finite.
Use the pole with positive imaginary part. Let \(m\) be the order
of the zero of \(M\) at \(z_*=\rho_*-\beta\). The transform of the
bounded complex density
\[
 J(v)=\mathbf1_{[0,L]}(v)\int_{[v,L]}
      \frac{(w-v)^{m-1}}{(m-1)!}e^{z_*(w-v)}\,d\nu(w)
\]
is \(M(z)/(z-z_*)^m\), by the integral Taylor remainder for
\(e^{zw}\) at \(z_*\) and \(M^{(j)}(z_*)=0\) for \(j<m\).
It takes a nonzero value at \(z_*\).
Thus \(\int J(v)g(u+v)\,dv\) equals a nonzero constant times
\(e^{z_*u}\), up to \(o(e^{(b_*-\beta)u})\), since the real
terms and those with real part at most \(\sigma\) have smaller
exponents. Termwise integration is justified by absolute convergence
on the fixed support. Also \(|J|\ll_* B\). The integral of
\(B|g|\) is therefore \(\gg_*e^{(b_*-\beta)u}\), whereas the
signed integral of \(Bg\) is \(o(e^{(b_*-\beta)u})\). Each sign
has a point of the required size there. Taking \(u=\log X+D+1\)
and \(C=\exp(2D+L+2)\) leaves unit logarithmic margins.

If \(\Theta>\beta\) is unattained, the same assertions hold for
every prescribed \(C>1\). As \(\sigma\uparrow\Theta\), the sum
defining \(L\) decreases to zero by dominated convergence at
\(\sigma=\beta\). Choose \(D>L\) and positive margins so small
that the resulting support with its margins fits in \([0,\log C]\).
Then choose \(b_*>\max(\sigma,\theta)\) as above.
\end{proof}

\begin{theorem}[A moment of the actual arithmetic tests]\label{scope:sparse-moment}
For any fixed \(p>0\), along any cofinal sequence of the specified
complete windows, or on \([Y,Z]=[X,CX]\) contained in those
windows with \(C\) from Lemma~\ref{scope:fixed-ratio} and \(X\to\infty\),
\[
 \lim_{Y\to\infty}
 \frac{1}{p\log Z}\log\!\left(
       \frac1{|\mathcal N_Y|}
       \sum_{m\in\mathcal N_Y}
                  \big(1+\max(Q(m^r),0)\big)^p\right)
       =\max(\Theta-\beta,0).
\]
The same limit holds using \(Q(x)=W_Z(x)/x^\beta\) with the ordinary
prime prefix of Proposition~\ref{pref:powers}; this version requires
only the fixed character. If \(\Theta>\beta\), then for every
\(a>0\) the proportion of indices with \(Q(m^r)>1\) exceeds
\(Z^{-a}\) in every sufficiently late such window.
\end{theorem}
\begin{proof}
For the geometric quantities let \(d(x)\) be the discrepancy
normalized by \(2/(|T|mb\delta_x)\), with \(m\) here the rank
parameter of the host, and let \(W(x)>0\) be the budget in the
same normalization. Then \(Q=d/W\). The discarded-twist bound,
the mean Fourier comparison and, for individual translates, the
overlap estimate give
\[
   |d(x)-F(x)|\ll_\psi x^{1/2+\varepsilon}.
\]
The page upper bound of Corollary~\ref{pref:scale} and the
\(A_3\) term in the budget in
Proposition~\ref{scope:reselection} give
\(2x^\beta/\delta_x\le W(x)\ll x^\beta/\delta_x\).
Here \(\delta_x^{-1}=x^{o(1)}\) uniformly. For the prime-prefix
version take \(d=W_Z,\ W=x^\beta\); the absolute error is
\(O(\sqrt{x})\). Thus in either case \(W\gg x^\beta\) and for
every fixed \(\xi>0\), \(W\ll_\xi x^{\beta+\xi}\).
The Hadamard bound \(|F(x)|\ll_\psi x^\Theta\) gives the upper
limit in the formula. When \(\Theta\le\beta\) the lower limit is
at least zero because each summand is at least one.

Suppose \(\Theta>\beta\). Fix \(0<\nu<\min(1,1/r)\). Choose
\(\max(\beta,\Theta-\nu/2)<\theta<\Theta\).
For all sufficiently large \(Z\), Lemma~\ref{scope:fixed-ratio}
gives \(t\asymp Z\) with \(F(t)>t^\theta\) in the interior of
\([Z/C,Z]\), a subinterval of either testing window.
The local variation bound in
Proposition~\ref{scope:fixed-powers} costs
\(O_\psi(t^{\Theta-\nu/2})=o(t^\theta)\) on
\([t-t^{1-\nu},t+t^{1-\nu}]\). This interval lies inside \([Y,Z]\)
eventually. Throughout it \(d(x)\gg t^\theta\), since
\(\theta>\beta>1/2+\varepsilon\), and hence
\(Q(x)\gg_\xi t^{\theta-\beta-\xi}\) for \(0<\xi<\theta-\beta\).
This lower bound holds simultaneously for every translate on the
overlap hosts.

Taking \(r\)-th roots of the two ends gives an interval of length
\(\asymp_r t^{1/r-\nu}\), which tends to infinity. It contains
\(\gg_r t^{1/r-\nu}\) integers, while
\(|\mathcal N_Y|\asymp_r Z^{1/r}\). The moment in parentheses is
therefore at least a constant times
\(Z^{-\nu+p(\theta-\beta-\xi)}\).
Let \(\nu,\xi\) tend to zero and \(\theta\) tend to
\(\Theta\) subject to the stated strict inequalities. This proves
the lower limit. For the final assertion each of those same indices
has \(Q>1\) eventually. Its proportion is bounded below by a
constant times \(Z^{-\nu}\), with \(\nu<a\) permissible.
\end{proof}

\begin{corollary}[Full estimates from tests with a sparse exceptional set]\label{scope:sparse-transfer}
Keep the common host hypotheses above. Suppose that on arbitrarily
late windows of either size all but \(O(Z^{-c}|\mathcal N_Y|)\) of the tested
powers satisfy \eqref{scope:average}, for some fixed \(c>0\).
Alternatively, on the overlap hosts it suffices that at the same
number of powers at least one actual word at one translate passes
the signed test of Corollary~\ref{scope:one-word}.
Then both \eqref{scope:average} and \eqref{scope:missing} hold for
all translates at every sufficiently late real scale of the
realized overlap family. The exponent and coefficient are the ones
in the two original displays. The same conclusion follows from
a subpower upper bound \(Z^{o(1)}\) for any one of the moments
in Theorem~\ref{scope:sparse-moment} on a cofinal subsequence.
\end{corollary}
\begin{proof}
A passing word gives \(D_u/E_x\le1\) at its translate by
\eqref{geom:finite-cancel} and the upper word loss. In this case
take \(Q=\min_u D_u/E_x\); for the mean take \(Q=\bar D/E_x\).
If \(\Theta>\beta\), the failing proportion exceeds \(Z^{-c/2}\)
eventually, contradicting the assumed estimate. The moment version
uses the positive limit when \(\Theta>\beta\).
Proposition~\ref{scope:fixed-powers} now applies at \(\Theta\le\beta\).
At every late scale on the overlap family it gives
\(\max_u D_u/E_x\le O_\psi(\delta_x)+o(1)\), which tends to zero
because \(4L_{w,x}|T_{2,x}|\le n_x/o\) and \(L_{w,x}\to\infty\).
In particular there is a uniform threshold for coefficient one.
\end{proof}

\begin{corollary}[Only the relative determinant in the word test]
\label{scope:signed-word-alone}
On the selected overlap kits satisfying Proposition~\ref{orb:selected},
use the local length and upper quotient comparison in
Section~\ref{geom:relative-finite-index} and the original size bounds.
For each positive-rank word on each translate one has
\[
        \frac{D_u}{E_x}\le\frac{P_{ui}}{E_x}
                  +O(x^{-4\varepsilon}+x^{-\beta}).
\]
Thus it suffices to bound \(P_{ui}/E_x\) by a fixed constant for
one actual word at one translate at all but a proportion \(O(Z^{-c})\)
of the powers in an arbitrarily late sequence of the specified
fixed-ratio intervals. These signed determinant tests imply both
original bounds on the full realized windows.
\end{corollary}
\begin{proof}
Put \(y=x^\varepsilon\), \(B=n_Tmb\), and use \(P\) of
Corollary~\ref{trmean:loss} for the nonnegative numerical bracket.
On the original ranks \(E_x\ge y^5 P\). Since the normalized
budget is \(W=2E_x/(B\delta_x)\ge 2x^\beta/\delta_x\), also
\(B/E_x\le x^{-\beta}\).
The local upper estimate is
\(D_u\le P_{ui}+(\mathcal A_u-G_u)+\ell_{ui}\).
Proposition~\ref{orb:selected} gives
\(\mathcal A_u-G_u\le CB,\ I_u\le Bx^{-8}\).
The level count bounds \(\ell_{ui}\le c_s I_u+K_7\) and
\(K_7\ll yP\) by the proof of Corollary~\ref{trmean:loss}.
Divide by \(E_x\) to obtain the displayed estimate. If
\(\Theta>\beta\), the lower-bound interval in the proof of
Theorem~\ref{scope:sparse-moment} has all translates above a fixed
positive power of \(Z\) and occupies more than a proportion
\(Z^{-c/2}\) of the tested powers for large \(Z\).
At least one of these powers lies outside the exceptional set,
contradicting its one-word bound. Apply Proposition~\ref{scope:fixed-powers}.
The original word loss is used only as an upper bound. The signed
relative determinant test on these intervals remains to be proved.
\end{proof}

\begin{corollary}[Intervals of actual points for bounded conductors]
\label{scope:bounded-conductor-intervals}
Fix a primitive target of conductor at most \(400\,000\), including
zeta, and fix \(0<\nu<1\). Take the same supplied hosts and realized
points with the uniform overlap and page estimates as above.
There is a finite \(C=C_{\psi,\varepsilon}>1\) such that every
sufficiently late \([X,CX]\) contained in one of their complete
windows contains a closed interval
\[
        [t-t^{1-\nu},\,t+t^{1-\nu}]
\]
on which \eqref{scope:average} and \eqref{scope:missing} both hold.
This conclusion uses no signed inequality on a word or translate.
\end{corollary}
\begin{proof}
Each primitive inducer of a power of \(\psi\) has conductor dividing
the conductor of \(\psi\). For the nonprincipal inducers the
finite-height verification of
\href{https://research-information.bris.ac.uk/files/67056136/platt_grh3.0.pdf}
{D.~J. Platt, \emph{Numerical computations concerning the GRH},
Theorem 9.1 in the author's manuscript} excludes real zeros in
\((1/2,1)\): such a zero would have ordinate zero. For \(\zeta\)
one can use the alternating series directly. If \(\Theta\le\beta\)
Proposition~\ref{scope:fixed-powers} covers every sufficiently
late tested scale and every intervening real scale on the hosts.
Otherwise choose \(\max(\beta,\Theta-\nu/2)<\theta<\Theta\).
Lemma~\ref{scope:fixed-ratio} supplies the negative excursion
in the interior of every specified \([X,CX]\).
The local variation bound makes \(F(x)<0\) throughout the stated
interval. Uniformly for the actual translates, the normalized
discrepancy is then at most \(O_\psi(x^{1/2+\varepsilon})\),
which their own normalized budget \(W(x)\gg x^\beta\) absorbs.
Both original estimates follow at the same points.
\end{proof}

\begin{corollary}[When a fixed fraction of a cell can pass]
\label{scope:relative-length}
Use the same bounded-conductor hypotheses and \(C\). In each such
cell \([X,CX]\), let \(H_X\) be the supremum of lengths of intervals
on which both estimates hold for every translate. If
\(\Theta\le\beta\), then \(H_X=(C-1)X\) eventually. If
\(\Theta>\beta\) and a retained zero has real part \(\Theta\), then
\(H_X\gg_{\psi,\varepsilon}X\) eventually. If \(\Theta>\beta\) and
the supremum is unattained, \(H_X/X\to0\) along the supplied cells.
\end{corollary}
\begin{proof}
The first case is Proposition~\ref{scope:fixed-powers}.
In the second case the zero attaining \(\Theta\) is nonreal, by the
real-zero exclusion in Corollary~\ref{scope:bounded-conductor-intervals}.
Use that zero as \(\rho_*\) in Lemma~\ref{scope:fixed-ratio}.
The isolated term has absolute size a constant times
\(e^{(\Theta-\beta)u}\), while the signed average against \(B\)
is \(o(e^{(\Theta-\beta)u})\). Thus at an interior point \(t\)
one has \(F(t)\le-c t^\Theta\) with \(c>0\) fixed. The variation
estimate \(\lvert F(t(1+v))-F(t)\rvert\ll_\psi t^\Theta\sqrt{|v|}\)
keeps \(F\) negative on a fixed relative neighborhood of \(t\).
Shrinking it to fit the logarithmic margins preserves length
\(\gg X\). The previous corollary's budget comparison applies there.

In the last case fix \(0<\kappa<C-1\). Any interval of length
\(\kappa X\) inside the cell contains a closed multiplicative
interval with ratio \(1+\kappa/(2C)\). The last assertion of
Lemma~\ref{scope:fixed-ratio} applies at this prescribed ratio.
For fixed \(\beta<\theta<\Theta\), it supplies a point \(t\) with
\(F(t)>t^\theta\) at all sufficiently large starting scales.
On the actual hosts the normalization error is
\(O_\psi(t^{1/2+\varepsilon})\) and the budget is at most
\(t^{\beta+o(1)}\), uniformly in the translate. Consequently
every translate exceeds its budget at that point. Letting
\(\kappa\) decrease to zero proves the limit.
\end{proof}
